\documentclass[12pt]{article}
\usepackage{geometry}
\geometry{ b4paper, total={220mm,320mm}, left=20mm, top=15mm,
headheight=33pt,includeheadfoot}

\usepackage[parfill]{parskip}  % Activate to begin paragraphs with an
% empty line rather than an indent

\usepackage{mathtools}
\usepackage{blindtext}
\usepackage{multicol}
\usepackage{listings, matlab-prettifier}
\usepackage{tikz}
\usepackage{wrapfig}
\usepackage{longtable}
\usepackage{booktabs, caption, colortbl}
\usepackage{pdflscape}
\usepackage{appendix}
\usepackage{titlesec}

\usepackage{pgfplots}
\pgfplotsset{compat=1.18}

\usepackage{amssymb}
\usepackage{enumitem}
\usepackage{amsmath}
%\usepackage{cleveref}

\usepackage{fancyhdr, lastpage, setspace}
\pagestyle{fancy}
\usepackage{caption, subcaption, zref-totpages}

\usepackage{fontspec}
\usepackage{microtype}

\usepackage{polyglossia}

\usepackage[dvipsnames]{xcolor} %red, green, blue, yellow, cyan,
% magenta, black, white

% \usepackage{multimedia}
% \usepackage{xmpmulti}
% \usepackage{media9}
% \usepackage{graphicx}
% \usepackage{animate}
%\def\pdfpageref{\pdffeedback pageref}
\usepackage{hyperref}
\hypersetup{ bookmarksopen=true, colorlinks=true, linkcolor=black,
    filecolor=magenta,
urlcolor=black, pdftitle={Markos Delaportas 1316 Thesis} }

% Load glossaries-extra with 'record' (triggers bib2gls) and 'abbreviations' support
% stylemods=long allows us to use table-based styles
\usepackage[record, abbreviations, stylemods={long}]{glossaries-extra}

% Set the style to a clean 2-column list
\setabbreviationstyle[acronym]{long-short} % First use: Long (Short)
\setglossarystyle{long} % A table style: Acronym | Description

% Load your external .bib file
\GlsXtrLoadResources[
  src={acronyms}, % The name of your .bib file without extension
  sort={en-US}    % Sort according to English rules
]


\setmainlanguage{english}
\setdefaultlanguage{english}
\setotherlanguages{greek}

\usepackage[ backend=biber, bibencoding=utf8, style=ieee]{biblatex}
\addbibresource{ref.bib}

\listfiles

\urlstyle{same}

\rhead{\includegraphics[width=1cm]{ece.png}}
\lhead{\includegraphics[height=1cm]{uowm.png}}

\fancyfoot{}
\fancyfoot[RO]{\thepage /\ztotpages}

\graphicspath{ {img} }
% \addmediapath{vdo}

%\setmainfont[Kerning=On,Ligatures=TeX]{Linux Libertine}
\setmainfont[Ligatures=TeX,Numbers=Lining]{Linux Libertine}
\setsansfont[Kerning=On]{Atkinson Hyperlegible}
\setmonofont[Scale=MatchLowercase]{Consolas}
%\setmathsf{Cambria-Math} \setmathtt{Consolas}

% Apply sans-serif font to all levels
\titleformat*{\section}{\Large\bfseries\sffamily}
\titleformat*{\subsection}{\large\bfseries\sffamily}
\titleformat*{\subsubsection}{\normalsize\bfseries\sffamily}

\renewcommand{\thesection}{\arabic{section}}
\renewcommand{\thesubsection}{\thesection.\alph{subsection}}
\renewcommand{\thesubsubsection}{\thesubsection.\roman{subsubsection}}

\newcommand{\ra}[1]{\renewcommand{\arraystretch}{#1}}

\onehalfspacing

\AddToHook{env/landscape/begin} {%
    \clearpage
    \pagestyle{empty}
    \AddToHook{shipout/background}[sven/page] {
        \put(0.9\paperwidth,-0.5\paperheight)%adapt values
    {\rotatebox{90}{\thepage}} }%
}

\AddToHook{env/landscape/after} {\RemoveFromHook{shipout/background}[sven/page]}

% Define a light grey background color
\definecolor{codebg}{rgb}{0.95,0.95,0.95}

% Configure the Listing Style
\lstset{
    style=Matlab-editor,           % Use standard MATLAB highlighting
    basicstyle=\ttfamily\small,    % Use the Consolas font defined above
    backgroundcolor=\color{codebg},% Grey background
    frame=none,                    % No black border (optional, remove if you want one)
    numbers=left,                  % Line numbers on the left
    numberstyle=\tiny\color{gray}, % Line numbers are small and grey
    numbersep=10pt,                % Space between line numbers and code
    xleftmargin=15pt,              % Indent the block to fit line numbers
    % This block forces numbers in the code to use the Math Font
    % and ensures they are "Lining Figures" (aligned on the baseline)
    literate=%
        {0}{{$0$}}1
        {1}{{$1$}}1
        {2}{{$2$}}1
        {3}{{$3$}}1
        {4}{{$4$}}1
        {5}{{$5$}}1
        {6}{{$6$}}1
        {7}{{$7$}}1
        {8}{{$8$}}1
        {9}{{$9$}}1
        {.}{{.}}1                  % Optional: make decimal points look like math too
}

\title{ \textsf{Over The Air Computing}\\
    \textsf{Master Thesis}\\
    \textsf{\Large Department of Electrical \& Computer Engineering}\\
    \textsf{\large University of Western Macedonia}
} \author{\textsf{Markos Delaportas} \footnote{E-mail: ece01316@uowm.gr}}
\date{\textsf{November 2025}}

\begin{document}

\maketitle

\textbf{ \textit{Abstract} -- Over The Air Computing (OTA Computing) is an
    emerging technique that leverages the superposition property of wireless
    channels to perform computations directly during the transmission process.
    This approach has the potential to significantly reduce latency and energy
    consumption in distributed systems, making it particularly suitable for
    applications in Internet of Things (IoT) networks, edge computing, and
    real-time data processing. This thesis explores the theoretical foundations
    of \gls{ota}, including channel modeling, signal processing techniques,
    and error analysis. Practical implementations and case
    studies demonstrating the effectiveness of OTA Computing in various
scenarios are also presented.}

{\let\clearpage\relax \tableofcontents} \thispagestyle{empty}
%\tableofcontents

% Add the Acronyms List here
\printunsrtglossary[type=abbreviations, title={List of Acronyms and Abbreviations}]
%\printglossary[type=\acronymtype, title=List of Abbreviations]
\clearpage

\section*{Introduction}
\addcontentsline{toc}{section}{\protect\numberline{}Introduction}
\subsection{Background and Motivation}
The increasing need for efficiency in computation-oriented applications
where the interest lies in a function of local information, rather than
the information itself \cite{sahin_survey_2023}. The conventional
approach of separating communication and computation often proves
suboptimal \cite{sahin_survey_2023,goldenbaum_analyzing_2011,
goldenbaum_robust_2013}. The potential of Over-the-Air Computation
(OAC) to achieve significantly higher computation rates by harnessing
interference through signal superposition in a wireless Multiple Access
Channel (MAC) \cite{sahin_survey_2023}.
\subsection{Problem Statement and Research Objectives}
Focusing on designing and validating a practical OAC network layout
suitable for implementation in a controllable
environment\cite{sahin_survey_2023}. The specific role of Software
Defined Radios (SDRs) and MATLAB/simulation tools for prototype
validation and addressing practical limitations such as synchronization
and channel impairments \cite{sahin_survey_2023,xiao_heterogeneous_2025,
evgenidis_waveform_2025}.
\subsection{Thesis Contributions and Organization}
\section{Over-the-Air Computation: Foundations and Techniques}
\subsection{Theoretical Foundations of OAC}
\begin{itemize}
    \item \textit{Signal Superposition in the Multiple Access Channel (MAC)}:
        The distinct feature of OAC where computation is handled by harnessing
        interference via simultaneous transmissions
        \cite{sahin_survey_2023,goldenbaum_robust_2013}.
    \item \textit{Nomographic Functions}: Defining the class of
        functions computable via
        OAC, which can be expressed as a post-processed sum of pre-processed
        sensor readings \cite{sahin_survey_2023,goldenbaum_analyzing_2011,
        goldenbaum_nomographic_2015,goldenbaum_robust_2013}.
    \item \textit{Examples of Nomographic Functions}: Arithmetic
        mean, weighted sum,
        maximum/minimum value, and modulo-2
        sum\cite{sahin_survey_2023,evgenidis_waveform_2025,goldenbaum_robust_2013}.
\end{itemize}
\subsection{OAC Schemes and Implementation Strategies}
\begin{itemize}
    \item \textit{Analog vs. Digital OAC}: Reviewing uncoded analog
        transmission (often
        optimal for simple Gaussian sensor networks) versus digital schemes
        designed for modern systems (e.g., using QAM or one-bit
        quantization)\cite{evgenidis_waveform_2025,zhu_one-bit_2021}.
    \item \textit{Coherent Schemes (CSI Dependent)}: Discussing
        methods requiring
        Channel State Information (CSI), such as Truncated Channel
        Inversion (TCI) or
        Zero-Forcing (ZF) coordination, for phase and amplitude
        alignment\cite{sahin_survey_2023,cao_optimized_2022,zhu_one-bit_2021}.
    \item \textit{Non-Coherent Schemes (CSI Independent)}: Examining
        techniques robust to
        synchronization errors, such as those using Orthogonal Signaling
        (FSK/PPM)\cite{sahin_survey_2023,yang_federated_2020}.
\end{itemize}
\subsection{Applications Context (e.g., Distributed Consensus)}
Application choice (e.g., Federated Edge Learning (FEEL) for gradient
    aggregation, or average consensus in (Wireless Sensor Networks (WSNs))
    
    \cite{sahin_survey_2023,yang_distributed_2025,yang_federated_2020,
    cao_optimized_2022,zhu_over--air_2024}.
    \section{System Model and Practical Challenges in OAC}
    \subsection{System Architecture and Channel Model}
    \begin{itemize}
        \item \textit{Modeling the Network}: Description of Edge Devices (EDs)
            transmitting to an Edge Server (ES)/Fusion Center (FC), possibly
            involving multiple antennas
            (\gls{mimo})\cite{sahin_survey_2023,yang_federated_2020,
            hellstrom_optimal_2023,cao_optimized_2022,zhu_over--air_2024}.
        \item \textit{Wireless Channel Model}: Consideration of
            channel impairments, including
            fading (e.g., Rayleigh fading), noise (\gls{awgn}), and path
            loss\cite{zhu_one-bit_2021,evgenidis_waveform_2025,cao_optimized_2022}.
    \end{itemize}
    \subsection{Impact of Practical Impairments on OAC Reliability}
    \subsubsection{Synchronization Errors}
    \textit{Time, Frequency, and Phase Offsets (TO, CFO, PO)}: The detrimental
    impact of imperfect time-of-arrivals and phase alignment on signal
    superposition and computation
    accuracy\cite{sahin_survey_2023,evgenidis_waveform_2025,hellstrom_optimal_2023}.

    \begin{lstlisting}
        % This is a comment
        x = linspace(0, 2*pi, 100);
        y = sin(x) + 0.5 * randn(size(x)); % Numbers like 0.5 will be in Math Font
        plot(x, y);
    \end{lstlisting}

    \subsubsection{Channel State Information (CSI) Imperfection}
    The need for accurate and fresh \gls{csi} at transmitters or
    receivers for many precoding/equalization techniques
    \cite{sahin_survey_2023,yang_federated_2020,evgenidis_waveform_2025}.
    \subsubsection{Waveform and Filter Design}
    Addressing time sampling error and intersymbol interference \gls{isi}
    inherent in digital communication transceivers when applied to OAC
    \cite{evgenidis_waveform_2025}.
    \subsection{Optimization Strategies to Enhance OAC Performance}
    \begin{itemize}
        \item \textit{Power Management}: Strategies to mitigate
            aggregation errors, such as
            optimizing power control policies to minimize Mean Squared Error
            \gls{mse} or maximize convergence speed
            \cite{sahin_survey_2023,cao_optimized_2022,evgenidis_waveform_2025}.
        \item \textit{Transceiver Design}: Joint design of
            pre-processing functions,
            precoders/beamformers, and post-processing functions (decoders) to
            match the channel to the desired function and minimize
            distortion \cite{goldenbaum_robust_2013,hellstrom_optimal_2023,
            sahin_survey_2023,yang_federated_2020}.
    \end{itemize}
    \section{Experimental Implementation using SDR and MATLAB}
    \subsection{Experimental Layout and Network Topology}
    \begin{itemize}
        \item \textit{Layout Description}: Presentation of the
            network structure (e.g., a
            single-cell OAC scenario or an ad-hoc network) using \gls{sdr}
            to represent
            Edge Devices (EDs) and a Fusion Center (FC)\cite{sahin_survey_2023}.
        \item \textit{Hardware and Software Tools}: Specific mention
            of the \gls{sdr} and the use of MATLAB for
            waveform generation, digital signal processing, control loops, and
            simulation\cite{sahin_survey_2023,xiao_heterogeneous_2025,
            yang_federated_2020,zhu_one-bit_2021}.
    \end{itemize}
    \subsection{Design and Implementation of the OAC Protocol}
    \subsubsection{Selection of the OAC Scheme for Demonstration}
    Choice of Target Function (e.g., Arithmetic Mean or Majority Vote):
    Justification based on practical implementation constraints of
    \gls{sdr}\cite{sahin_survey_2023}.
    \subsubsection{Synchronization Mechanisms}
    Detailed methodology for achieving coarse block synchronization or
    mitigating synchronization errors (TO/CFO) in the \gls{sdr}/MATLAB
    environment (e.g., based on trigger signals or custom control
    loops)\cite{sahin_survey_2023,zhu_one-bit_2021,yang_federated_2020}.
    \subsubsection{Pre-processing and Channel Compensation Implementation}
    Implementation of the chosen modulation/encoding scheme and any
    necessary pre-equalization (e.g., Channel Inversion precoding if a
        coherent scheme is selected, or non-coherent coding if robustness is
    prioritized)\cite{sahin_survey_2023,zhu_one-bit_2021}.
    \subsection{Measurement and Data Collection Methodology}
    Procedure for varying network parameters (e.g., number of transmitting
    devices, transmit SNR)\cite{evgenidis_waveform_2025}.
    \section{Experimental Results and Discussion}
    \subsection{Validation of the Implemented OAC Network}
    Demonstration of successful signal superposition and computation via the
    SDR platform\cite{sahin_survey_2023,goldenbaum_analyzing_2011}.
    Comparison of achieved computation result against the ideal theoretical
    value\cite{sahin_survey_2023,zhu_one-bit_2021}.
    \subsection{Performance Evaluation under Impairments}
    Quantifying the performance metrics, such as Mean Squared Error (MSE) or
    Computation Error Rate (CER), of the implemented
    scheme\cite{sahin_survey_2023,evgenidis_waveform_2025}. Analysis of the
    impact of deliberately introduced synchronization errors or varying
    channel noise (SNR) on the computation
    accuracy\cite{sahin_survey_2023,evgenidis_waveform_2025}. Comparison
    with Separation Schemes (Simulated Baseline); Illustrating the resource
    savings and computation rate achieved by the OAC approach relative to
    traditional orthogonal access (e.g.,
    \gls{tdma})\cite{goldenbaum_robust_2013,sahin_survey_2023,goldenbaum_analyzing_2011}.
    \section{Conclusion and Future Work}
    \subsection{Summary of Thesis Findings}
    Recap of the feasibility and performance demonstrated by the SDR-based
    OAC network.
    \subsection{Future Research Directions}
    Exploring more complex OAC architectures, such as multi-antenna systems
    \gls{mimo} or using Reconfigurable Intelligent Surfaces
    (RIS)\cite{evgenidis_waveform_2025,zhu_over--air_2024}. Further
    investigation into standardized OAC protocols for compatibility with
    modern wireless networks (e.g., LTE/5G
    protocols)\cite{zhu_one-bit_2021,sahin_survey_2023}.

    \printbibliography

    \end{document}
